% Options for packages loaded elsewhere
% Options for packages loaded elsewhere
\PassOptionsToPackage{unicode}{hyperref}
\PassOptionsToPackage{hyphens}{url}
\PassOptionsToPackage{dvipsnames,svgnames,x11names}{xcolor}
%
\documentclass[
  a4paperpaper,
  DIV=11,
  numbers=noendperiod]{scrartcl}
\usepackage{xcolor}
\usepackage{amsmath,amssymb}
\setcounter{secnumdepth}{2}
\usepackage{iftex}
\ifPDFTeX
  \usepackage[T1]{fontenc}
  \usepackage[utf8]{inputenc}
  \usepackage{textcomp} % provide euro and other symbols
\else % if luatex or xetex
  \usepackage{unicode-math} % this also loads fontspec
  \defaultfontfeatures{Scale=MatchLowercase}
  \defaultfontfeatures[\rmfamily]{Ligatures=TeX,Scale=1}
\fi
\usepackage{lmodern}
\ifPDFTeX\else
  % xetex/luatex font selection
\fi
% Use upquote if available, for straight quotes in verbatim environments
\IfFileExists{upquote.sty}{\usepackage{upquote}}{}
\IfFileExists{microtype.sty}{% use microtype if available
  \usepackage[]{microtype}
  \UseMicrotypeSet[protrusion]{basicmath} % disable protrusion for tt fonts
}{}
\makeatletter
\@ifundefined{KOMAClassName}{% if non-KOMA class
  \IfFileExists{parskip.sty}{%
    \usepackage{parskip}
  }{% else
    \setlength{\parindent}{0pt}
    \setlength{\parskip}{6pt plus 2pt minus 1pt}}
}{% if KOMA class
  \KOMAoptions{parskip=half}}
\makeatother
% Make \paragraph and \subparagraph free-standing
\makeatletter
\ifx\paragraph\undefined\else
  \let\oldparagraph\paragraph
  \renewcommand{\paragraph}{
    \@ifstar
      \xxxParagraphStar
      \xxxParagraphNoStar
  }
  \newcommand{\xxxParagraphStar}[1]{\oldparagraph*{#1}\mbox{}}
  \newcommand{\xxxParagraphNoStar}[1]{\oldparagraph{#1}\mbox{}}
\fi
\ifx\subparagraph\undefined\else
  \let\oldsubparagraph\subparagraph
  \renewcommand{\subparagraph}{
    \@ifstar
      \xxxSubParagraphStar
      \xxxSubParagraphNoStar
  }
  \newcommand{\xxxSubParagraphStar}[1]{\oldsubparagraph*{#1}\mbox{}}
  \newcommand{\xxxSubParagraphNoStar}[1]{\oldsubparagraph{#1}\mbox{}}
\fi
\makeatother


\usepackage{longtable,booktabs,array}
\usepackage{calc} % for calculating minipage widths
% Correct order of tables after \paragraph or \subparagraph
\usepackage{etoolbox}
\makeatletter
\patchcmd\longtable{\par}{\if@noskipsec\mbox{}\fi\par}{}{}
\makeatother
% Allow footnotes in longtable head/foot
\IfFileExists{footnotehyper.sty}{\usepackage{footnotehyper}}{\usepackage{footnote}}
\makesavenoteenv{longtable}
\usepackage{graphicx}
\makeatletter
\newsavebox\pandoc@box
\newcommand*\pandocbounded[1]{% scales image to fit in text height/width
  \sbox\pandoc@box{#1}%
  \Gscale@div\@tempa{\textheight}{\dimexpr\ht\pandoc@box+\dp\pandoc@box\relax}%
  \Gscale@div\@tempb{\linewidth}{\wd\pandoc@box}%
  \ifdim\@tempb\p@<\@tempa\p@\let\@tempa\@tempb\fi% select the smaller of both
  \ifdim\@tempa\p@<\p@\scalebox{\@tempa}{\usebox\pandoc@box}%
  \else\usebox{\pandoc@box}%
  \fi%
}
% Set default figure placement to htbp
\def\fps@figure{htbp}
\makeatother





\setlength{\emergencystretch}{3em} % prevent overfull lines

\providecommand{\tightlist}{%
  \setlength{\itemsep}{0pt}\setlength{\parskip}{0pt}}



 


\newcommand{\B}[1]{\boldsymbol{#1}}
\newcommand{\+}{\phantom{-}}
\newcommand{\op}[1]{\operatorname{#1}}
\KOMAoption{captions}{tableheading}
\makeatletter
\@ifpackageloaded{caption}{}{\usepackage{caption}}
\AtBeginDocument{%
\ifdefined\contentsname
  \renewcommand*\contentsname{Table of contents}
\else
  \newcommand\contentsname{Table of contents}
\fi
\ifdefined\listfigurename
  \renewcommand*\listfigurename{List of Figures}
\else
  \newcommand\listfigurename{List of Figures}
\fi
\ifdefined\listtablename
  \renewcommand*\listtablename{List of Tables}
\else
  \newcommand\listtablename{List of Tables}
\fi
\ifdefined\figurename
  \renewcommand*\figurename{Figure}
\else
  \newcommand\figurename{Figure}
\fi
\ifdefined\tablename
  \renewcommand*\tablename{Table}
\else
  \newcommand\tablename{Table}
\fi
}
\@ifpackageloaded{float}{}{\usepackage{float}}
\floatstyle{ruled}
\@ifundefined{c@chapter}{\newfloat{codelisting}{h}{lop}}{\newfloat{codelisting}{h}{lop}[chapter]}
\floatname{codelisting}{Listing}
\newcommand*\listoflistings{\listof{codelisting}{List of Listings}}
\makeatother
\makeatletter
\makeatother
\makeatletter
\@ifpackageloaded{caption}{}{\usepackage{caption}}
\@ifpackageloaded{subcaption}{}{\usepackage{subcaption}}
\makeatother
\usepackage{bookmark}
\IfFileExists{xurl.sty}{\usepackage{xurl}}{} % add URL line breaks if available
\urlstyle{same}
\hypersetup{
  pdftitle={Analysis of Evolutionary Algorithms for Optimisation},
  pdfauthor={Shvet Maharaj},
  colorlinks=true,
  linkcolor={blue},
  filecolor={Maroon},
  citecolor={Blue},
  urlcolor={Blue},
  pdfcreator={LaTeX via pandoc}}


\title{Analysis of Evolutionary Algorithms for Optimisation}
\author{Shvet Maharaj}
\date{12 January 2026}
\begin{document}
\maketitle

\renewcommand*\contentsname{Table of contents}
{
\hypersetup{linkcolor=}
\setcounter{tocdepth}{4}
\tableofcontents
}

\section{Design}\label{design}

\subsection{Update procedures}\label{update-procedures}

Where noted, a \emph{strictly improving update procedure} will be used
alongside the updating equations, defined as:
\[\B x_{k+1} = \begin{cases}\B v_{k+1}&\text{if} \quad \op f(\B v_{k+1})< \op f(\B x_k),\\\B x_k & \text{otherwise.}\end{cases}\]
Where \(\B x_k\) is the best performing point at iteration \(k\), and
\(\B v_{k+1}\) is the candidate point for the next iteration, \(k+1\).
This candidate point is what is generated by the updating equations.

Alternatively, the update procedure could be a simple iterative updating
procedure, where each updated point, \(\B v_{k+1}\), is used directly by
the algorithm in the next iteration (thus, the algorithm updates on
\(\B v_k\), instead of \(\B x_k\)). This would correspond more so to the
deterministic algorithms. The best solution, \(\B x_k\) is simply stored
as the algorithm progresses.

\subsection{Updating equations}\label{updating-equations}

At the core of this exploration of the genetic algorithm, and other
methods for optimisation, are the \emph{updating equations}. These
define how exactly the optimisation paradigm attempts to iteratively
improve a candidate solution, \(\B v\).

\subsubsection{Hessian methods}\label{hessian-methods}

\paragraph{Newton-Raphson}\label{newton-raphson}

\[\B x_{k+1} = \B v_k -\frac{f'\left(\B v_k\right)}{f''\left(\B v_k\right)}\]

\subsubsection{Gradient methods}\label{gradient-methods}

\paragraph{Gradient descent}\label{gradient-descent}

\[\B v_{k+1} = \B v_k -\nu \nabla f\left( \B v_k \right),\] where
\(\nu\) is the learning rate.

\subsubsection{Function value methods}\label{function-value-methods}

Also known as Direct Search, Pattern Search, or Black-Box Methods. These
algorithms do not make use of gradient nor of hessian information. This
makes them far easier to apply to any context.

\paragraph{Random search (RS)}\label{random-search-rs}

The specific version of the random search algorithm is informed by the
descriptions given by (\textbf{zabinsky2009?}). Algorithms are described
by being `global' or `local'.

\textbf{Pure Random Search (PRS)} repeatedly performs a global search
with a simply improving update procedure.

The updating equation is defined as, \[\B v_{k+1} = \B u.\]

It is notable that neither of the terms \(\B x_k\) nor \(\B v_k\) ---
the best or candidate points found by or for iteration \(k\) --- are
present in this formula.

Each generated point is \emph{independent} of any previous point. The
major benefit of this is PRS is `embarrassingly parallelisable'. This is
the only optimisation algorithm that will have this property in this
paper.

We test also an additional variant of RS, where the global search is
either performed solely for the starting point), \(\B x_0\), or skipped
entirely in favour of a user-defined starting point. Subsequent points
are generated locally from the hypersphere (the `(\(N-1\))-sphere', or
the `surface of the \(N\)-ball') around \(\B x_k\).

This is Rastrigin's \textbf{Fixed Step Size Random Search
(FSSRS)}(\textbf{rastrigin1963?}).

An efficient way to sample uniformly from the surface of an \(N\)-ball
is given by Marsaglia(\textbf{marsaglia1972?}).

\[\B v_{k+1} = \B x_k +h \B u,\]

Here, \(\B u\) is a normalised \(N\)-dimensional vector sample from the
standard normal distribution such that,

\(\B u = \frac1{\lVert \B z \rVert}\B z\) and
\(\B z \sim \mathcal{N}_N\left(0,1\right),\) and where \(h\) is a step
size hyperparameter.

\paragraph{Genetic algorithms}\label{genetic-algorithms}

Genetic algorithms are a subset of evolutionary algorithms, typified by
the use of all of three genetic operators: crossover, selection, and
mutation. A \emph{real-coded} (in reference to \(\mathbb R\), the real
number set) genetic algorithm is used for greater applicability to
arbitrary problems and models, despite genetic processes occurring with
a discrete set of genes (ACGT) in nature.

Notation does shift here, due to the genetic algorithm necessarily
requiring multiple candidates for each iteration. Our previous notation
referencing a single candidate per iteration is thus no longer feasible.

We are implementing genetic algorithms using the \texttt{GA} package, as
described by Scrucca(\textbf{scrucca2013?}).

A brief explanation of the algorithm follows:

\[\B \Theta^{\left(k+1\right)} = \Theta^{\left(k\right)}\] \#\#\#\#\#
Fitness evaluation Fitness for each of the `genotypes' (candidate
solutions), \(\theta_i^{\left(k\right)}\) is calculated, where \(i\) is
an arbitrary index of the genotype in the generation \(k\). This is
typically the objective function of interest,
\(\op f\left(\cdot\right)\).

\subparagraph{Selection}\label{selection}

A subset of genotypes from the current generation's population becomes
selected as a separate population known as the `reproducing population'
Genotypes are selected (with replacement) through a weighted selection
process, based typically proportionally on the fitness of each genotype,
\(\op f\left(\theta_i^{\left(k\right)}\right)\), with probability
\(p_i^{\left(k\right)}\).

An elitist strategy can be used, which means that even if the best
performers are unluckily not selected, they still make it to the next
population. By default \texttt{gareal\_lsSelection} is used, which is
proportional selection with fitness linear scaling.

\subparagraph{Crossover}\label{crossover}

The core of a genetic algorithm, the crossover operation, brings a means
to generate additional `child' (next generation) genotypes through parts
of the `genetic information' (parameter values) from either of the
`parent' (current generation) genotypes. Alongside mutation, crossover
applied on the reproducing population creates the new, next generation
population. By default, \texttt{gareal\_laCrossover} is used, which is
local arithmetic crossover.

\subparagraph{Mutation}\label{mutation}

The randomly shifting of the values of `genes' (parameters) within
genotypes. The probability of mutation can be set by the user. Alongside
crossover, mutation applied on the reproducing population creates the
new, next generation population. By default, \texttt{gareal\_raMutation}
is used, which is uniform random mutation.

Crossover and mutation operations increase diversity in a population,
allowing us to explore more of the sample space (exploration), while
selection operations decrease diversity in favour of increased fitness
(exploitation).




\end{document}
